\chapter{Fatty Liver}

\section*{Definition and Overview}
Fatty liver, or hepatic steatosis, is the accumulation of excess fat within liver cells. A small amount of fat in the liver is normal, but when more than 5\% of the liver's weight is fat, it is called fatty liver disease. The most common form, non-alcoholic fatty liver disease (NAFLD), occurs in people who drink little or no alcohol and is closely linked to being overweight, having diabetes, and having high cholesterol. A less common form, alcoholic fatty liver, is caused by heavy drinking. Fatty liver itself usually causes no symptoms, but in some people it progresses to inflammation, scarring (fibrosis), cirrhosis, and liver failure, and it increases the risk of liver cancer. Because it is reversible in its early stages, recognising and addressing it is one of the most important opportunities in preventive medicine.

\section*{Epidemiology}
Fatty liver has become the most common liver disease in Australia and worldwide, closely tracking the rise in obesity and type 2 diabetes. Estimates suggest that around one in three Australian adults has fatty liver, and the proportion rises to well over half of people who are overweight or have diabetes. It also affects a growing number of children and adolescents. Only a minority of people with fatty liver develop significant liver scarring, but because so many people are affected, it is now one of the leading causes of cirrhosis and of liver transplants in many countries. Alcoholic fatty liver affects a smaller group but progresses faster when drinking continues.

\section*{Aetiology and Pathophysiology}
\begin{itemize}
    \item \textbf{Energy excess:} consuming more energy than the body uses causes fat to be stored not only under the skin but inside organs, including the liver
    \item \textbf{Insulin resistance:} the body's reduced response to insulin promotes release of fatty acids into the blood and increases fat production inside liver cells
    \item \textbf{Metabolic syndrome:} obesity, diabetes, high blood pressure, and abnormal cholesterol commonly coexist with fatty liver and worsen it
    \item \textbf{Alcohol:} alcohol is directly toxic to liver cells and stimulates fat deposition, producing alcoholic fatty liver
    \item \textbf{Other causes:} some medicines (such as steroids and certain chemotherapy drugs), hepatitis C, rapid weight loss, and rare genetic conditions can cause fat accumulation
    \item \textbf{Progression:} fat accumulation can trigger inflammation (steatohepatitis), which drives scarring; once significant scar tissue forms, the liver cannot repair itself fully
    \item \textbf{Gut and genetic factors:} the gut microbiome and inherited variations, such as in the PNPLA3 gene, influence who develops progressive disease
\end{itemize}

\section*{Clinical Features}
Fatty liver is usually silent, which is why it is often found by chance.
\begin{itemize}
    \item No symptoms in the vast majority of people, who are diagnosed incidentally on blood tests or ultrasound
    \item A dull ache or fullness in the right upper abdomen in some people with an enlarged liver
    \item Fatigue and a general sense of being unwell in a minority
    \item Features of the metabolic syndrome: weight gain, high blood sugar, and high blood pressure
    \item With advanced disease: jaundice, abdominal swelling, easy bruising, and confusion (encephalopathy)
    \item Alcoholic fatty liver may be accompanied by the effects of alcohol on other organs
\end{itemize}

\section*{Differential Diagnosis}
The same pattern of liver test abnormalities can arise from many conditions.
\begin{itemize}
    \item Viral hepatitis, including hepatitis B and C
    \item Alcoholic liver disease, distinguished mainly by the drinking history
    \item Autoimmune hepatitis and other immune-mediated liver diseases
    \item Drug-induced liver injury from prescription or herbal medicines
    \item Haemochromatosis (iron overload) and Wilson disease (copper overload)
    \item Biliary disease, including primary biliary cholangitis
\end{itemize}

\section*{When to Seek Medical Care}
Anyone with risk factors such as obesity, diabetes, or heavy drinking should have their liver checked by a doctor, especially if they feel persistently tired or have discomfort in the upper abdomen. Liver blood tests are part of routine health checks and should be reviewed with a doctor.

\textbf{Call 000 (or local emergency services) for:}
\begin{itemize}
    \item Vomiting blood or passing black, tarry stools, which may indicate bleeding from varices in advanced disease
    \item Sudden confusion, drowsiness, or personality change in someone with known liver disease
    \item Severe abdominal pain, which could indicate a complication such as infection of abdominal fluid
    \item Sudden swelling of the abdomen or legs with breathlessness
\end{itemize}

\section*{Diagnosis and Investigations}
\begin{itemize}
    \item \textbf{Liver function tests:} elevated liver enzymes (ALT and AST) raise suspicion, though levels can be normal in fatty liver
    \item \textbf{Abdominal ultrasound:} the usual first imaging test, showing the typical bright, fatty liver appearance
    \item \textbf{Metabolic blood tests:} fasting glucose, HbA1c, cholesterol, and triglycerides to assess coexisting metabolic disease
    \item \textbf{Exclusion of other causes:} hepatitis serology, iron studies, autoimmune markers, and alcohol history
    \item \textbf{Fibrosis assessment:} specialised blood scores (such as FIB-4) and elastography (FibroScan) measure liver stiffness and scarring without a biopsy
    \item \textbf{Liver biopsy:} reserved for uncertain or severe cases to grade inflammation and fibrosis
\end{itemize}

\section*{Management}
\begin{itemize}
    \item \textbf{Weight loss:} the cornerstone of treatment; losing 5--10\% of body weight reduces liver fat significantly, and greater loss can reverse inflammation and scarring
    \item \textbf{Diet:} a balanced diet with reduced processed foods, sugary drinks, and saturated fat, modelled on the Mediterranean diet
    \item \textbf{Physical activity:} regular exercise reduces liver fat even without weight loss
    \item \textbf{Diabetes and cholesterol control:} good glycaemic control and lipid management slow progression
    \item \textbf{Alcohol:} avoid or strictly limit alcohol; people with alcoholic fatty liver should stop completely
    \item \textbf{Medication review:} stop or substitute medicines that contribute to liver fat where possible
    \item \textbf{Vitamin E and other treatments:} selected people may benefit from vitamin E or newer medicines such as resmetirom, under specialist advice
    \item \textbf{Monitoring:} regular liver tests and fibrosis measurement to track progress
\end{itemize}

\section*{Complications}
\begin{itemize}
    \item Steatohepatitis, the inflammatory form that drives progression
    \item Fibrosis and cirrhosis, with the complications of liver failure, ascites, and variceal bleeding
    \item Hepatocellular carcinoma (liver cancer), which can develop even without cirrhosis in fatty liver
    \item Increased cardiovascular risk: heart attack and stroke are more common in people with fatty liver than liver failure is
    \item End-stage liver disease requiring transplantation
\end{itemize}

\section*{Prognosis and Outcomes}
The outlook depends on how much scarring is present. People with simple fatty liver who achieve weight loss and lifestyle change very often normalise their liver tests and liver fat within months, and even early fibrosis can regress. Once cirrhosis is established, the liver damage is permanent, although further harm can be slowed. The greatest threat to most people with fatty liver is not the liver itself but heart disease, which is why cardiovascular risk must be managed as carefully as the liver. With sustained weight loss, the risk of progression and complications falls substantially.

\section*{Prevention and Self-Care}
\begin{itemize}
    \item Maintain a healthy weight and prevent excess weight gain through diet and activity
    \item Limit sugary drinks, highly processed foods, and saturated fats
    \item Include vegetables, fruit, wholegrains, legumes, fish, and olive oil in the diet
    \item Exercise most days, building to at least 150 minutes of moderate activity a week
    \item Limit alcohol to well below the national guidelines, or stop completely if the liver is affected
    \item Attend regular health checks, including blood pressure, blood sugar, and liver tests
    \item Ask about a FibroScan if you have risk factors for fatty liver
\end{itemize}

\section*{Sources and Further Reading}
The following reputable sources provide further information for healthcare students and patients seeking reliable, current advice about fatty liver.
\begin{itemize}
    \item Healthdirect Australia. \textit{Fatty liver disease.} https://www.healthdirect.gov.au/fatty-liver
    \item Gastroenterology Society of Australia (GESA). \textit{Fatty liver patient information.} https://www.gesa.org.au/
    \item Diabetes Australia. \textit{Fatty liver and diabetes.} https://www.diabetesaustralia.com.au/
    \item NHS. \textit{Non-alcoholic fatty liver disease.} https://www.nhs.uk/conditions/non-alcoholic-fatty-liver-disease/
    \item Mayo Clinic. \textit{NAFLD.} https://www.mayoclinic.org/
    \item World Health Organization. \textit{Nutrition and liver health.} https://www.who.int/
\end{itemize}
